\textbf{Proof.}
Fix \(\mathcal A\in\mathfrak H_\alpha\), and take the \(M_n=m_n+1\) hypotheses from the many-bump lemma.  Write
\[
 d_n=\min_{j\ne\ell}\|\tau_j-\tau_\ell\|_\infty\ge c_dr_n,
 \qquad
 D_n(\mathcal A)=\sup_{v\in\mathcal I}\{U_n(v)-L_n(v)\}. \tag{1}
\]
Let \(P_j^n\) be the sample law under \(\tau_j\) and exact code zero, and let \(\mathcal C_{j,n}\) be the event that the entire curve \(\tau_j\) lies in the band.  Uniform asymptotic honesty gives
\[
 \kappa_n:=\min_{0\le j\le m_n}P_j^n(\mathcal C_{j,n})
 \ge1-\alpha-o(1). \tag{2}
\]

For a manifestly measurable decoder, evaluate all candidate curves and the band at the finite set of bump centers.  Let \(\widehat J\) be the unique index whose resulting finite vector is contained coordinatewise in those band intervals, with an arbitrary value if uniqueness fails.  Full-curve containment implies finite-vector containment.  If \(\mathcal C_{j,n}\) occurs and \(D_n(\mathcal A)<d_n\), no second curve can be contained: two candidates differ by at least \(d_n\) at one of the selected centers.  Markov's inequality therefore yields
\[
 \frac1{M_n}\sum_{j=0}^{m_n}P_j^n(\widehat J=j)
 \ge \kappa_n-
 \frac1{M_nd_n}\sum_{j=0}^{m_n}E_jD_n(\mathcal A). \tag{3}
\]

We derive the testing inequality used to bound the left side.  Let \(J\) be uniform on the hypotheses and draw the observations from \(P_J^n\).  For \(\overline P=M_n^{-1}\sum_jP_j^n\), the information-centroid identity and nonnegativity of divergence give
\[
 I(J;O_1^n)=\frac1{M_n}\sum_j\operatorname{KL}(P_j^n,\overline P)
 \le\frac1{M_n}\sum_j\operatorname{KL}(P_j^n,P_0^n)
 \le c_0\log m_n. \tag{4}
\]
If \(p_e=P(\widehat J\ne J)\), conditioning first on the error indicator gives
\[
 H(J\mid O_1^n)\le\log2+p_e\log M_n.
\]
Since \(H(J)=\log M_n\), equation (4) implies
\[
 \frac1{M_n}\sum_jP_j^n(\widehat J=j)
 \le\frac{c_0\log m_n+\log2}{\log M_n}. \tag{5}
\]
Combining (2), (3), and (5), then choosing the bump amplitude so that \(c_0<(1-\alpha)/2\), gives
\[
 \max_{0\le j\le m_n}E_jD_n(\mathcal A)
 \ge d_n\{1-\alpha-c_0-o(1)\}
 \ge c r_n. \tag{6}
\]
Every pair in (6) belongs to \(\mathcal Q_{\mathrm{spec},n}\); hence (6) is bounded above by the supremum risk in the theorem.  Taking the lower limit proves the all-band lower bound.  The honest-spectral-band theorem supplies the matching \(Cr_n\) upper bound for the particular spectral band on the same pair class.  No symmetry, fixed-width, or contour restriction entered (1)--(6), which proves the minimax conclusion. \(\square\)